\documentclass[11pt]{article}

% Change "review" to "final" to generate the final (sometimes called camera-ready) version.
% Change to "preprint" to generate a non-anonymous version with page numbers.
\usepackage[review]{acl}

% Standard package includes
\usepackage{times}
\usepackage{latexsym}

% For proper rendering and hyphenation of words containing Latin characters (including in bib files)
\usepackage[T1]{fontenc}

% This assumes your files are encoded as UTF8
\usepackage[utf8]{inputenc}

% This is not strictly necessary, and may be commented out,
% but it will improve the layout of the manuscript,
% and will typically save some space.
\usepackage{microtype}

% This is also not strictly necessary, and may be commented out.
% However, it will improve the aesthetics of text in
% the typewriter font.
\usepackage{inconsolata}

% Including images in your LaTeX document requires adding additional package(s)
\usepackage{graphicx}

% Additional packages for this paper
\usepackage{amsmath}
\usepackage{amssymb}
\usepackage{booktabs}
\usepackage{xcolor}
\usepackage{listings}
\usepackage{algorithm}
\usepackage{algorithmic}
\usepackage{multirow}
\usepackage{url}
\usepackage{tabularx}

% Code listing setup
\definecolor{codegreen}{rgb}{0,0.6,0}
\definecolor{codegray}{rgb}{0.5,0.5,0.5}
\definecolor{codepurple}{rgb}{0.58,0,0.82}
\definecolor{backcolour}{rgb}{0.97,0.97,0.95}

\lstdefinestyle{mystyle}{
    backgroundcolor=\color{backcolour},
    commentstyle=\color{codegreen},
    keywordstyle=\color{magenta},
    numberstyle=\tiny\color{codegray},
    stringstyle=\color{codepurple},
    basicstyle=\ttfamily\scriptsize,
    breakatwhitespace=false,
    breaklines=true,
    captionpos=b,
    keepspaces=true,
    numbers=none,
    showspaces=false,
    showstringspaces=false,
    showtabs=false,
    tabsize=2,
    frame=single,
}

\lstset{style=mystyle}

\title{AICC-IntelliDoc: A Multi-Agent Orchestration Framework for Document-Intelligent Conversational AI}

\author{First Author \\
  Affiliation / Address line 1 \\
  Affiliation / Address line 2 \\
  \texttt{email@domain} \\\And
  Second Author \\
  Affiliation / Address line 1 \\
  Affiliation / Address line 2 \\
  \texttt{email@domain} \\}

\begin{document}
\maketitle

\begin{abstract}
Translating visual AI agentic workflows into production-ready chatbots remains a critical challenge. Existing solutions fragment across multiple dimensions: multi-agent frameworks excel at coordination but lack deployment capabilities; visual workflow tools offer interfaces but provide workflow-level RAG configuration; RAG systems require manual integration. This paper introduces \textbf{AICC-IntelliDoc}, the first open-source framework providing a complete pipeline from visual multi-agent workflow design to deployable chatbot. We present three main contributions: (1) \textbf{Dual-Mode Human-in-the-Loop Architecture} supporting both Admin-in-the-Loop and User-in-the-Loop modes for executive oversight and end-user interactions, (2) \textbf{DocAware: Unified RAG with Hierarchical Document Management} enabling per-agent document-aware retrieval with folder and file-level content filtering, and (3) \textbf{Deployment Architecture: Workflow-to-Chatbot Translation} providing HTML embed code generation with allowed origins management and per-origin rate limiting. The framework implements five-layer project isolation ensuring complete multi-tenant separation. We present system architecture and implementation details demonstrating these unique capabilities.
\end{abstract}

\section{Introduction}

The evolution of foundational language models has catalyzed a paradigm shift in how we build intelligent applications. However, translating AI agentic workflows into production-ready chatbots remains a significant challenge. Existing solutions fragment across multiple dimensions: multi-agent frameworks excel at agent coordination but lack built-in multi-tenant isolation and deployment capabilities; visual workflow tools offer production-ready interfaces but provide workflow-level RAG configuration rather than per-agent document awareness; RAG systems require manual integration and lack hierarchical document management. This fragmentation creates a critical gap: organizations cannot easily translate visual workflow designs into deployable chatbots.

\textbf{AICC-IntelliDoc} addresses this gap by providing a unified framework that translates visual AI agentic workflows into production-ready chatbots. The framework's core strength is its ability to bridge the entire pipeline from visual workflow design to deployable chatbot with HTML embed code integration. This paper makes three main contributions: (1) \textbf{Dual-Mode Human-in-the-Loop Architecture}—UserProxyAgent supporting both Admin-in-the-Loop and User-in-the-Loop, enabling executive oversight and end-user interactions in deployed chatbots, (2) \textbf{DocAware: Unified RAG with Hierarchical Document Management}—per-agent RAG configuration with hierarchical document management enabling folder and file-level content filtering, and (3) \textbf{Deployment Architecture: Workflow-to-Chatbot Translation}—HTML embed code generation with allowed origins management and per-origin rate limiting for seamless website integration. The framework implements five-layer project isolation ensuring complete multi-tenant separation.

\section{Related Work and Gap Analysis}

\textit{[This section to be completed by authors. Will address multi-agent frameworks, visual workflow tools, RAG integration approaches, and HITL implementations with accurate comparative analysis.]}

The multi-agent ecosystem has fragmented into specialized tools—multi-agent frameworks excel at coordination, RAG systems dominate retrieval, and visual builders democratize access—but no single solution unifies these capabilities with enterprise-grade multi-tenancy. This section identifies critical gaps that AICC-IntelliDoc addresses.

\textbf{Multi-Tenant Isolation Gap:} Existing multi-agent frameworks lack built-in multi-tenant isolation, requiring custom implementation with separate runtimes per tenant or providing only thread-based or project-level separation. This gap is significant for enterprise deployments where customer data must be securely separated. AICC-IntelliDoc addresses this with a five-layer isolation architecture providing physical separation at the vector database level combined with logical separation at application layers.

\textbf{Visual-to-Production Deployment Gap:} Visual workflow tools and code-first frameworks serve different audiences with no bridging mechanism to production chatbot deployment. Existing solutions require custom deployment infrastructure or lack HTML embed code generation and allowed origins management for seamless website integration. AICC-IntelliDoc bridges this divide by translating visual workflow designs into deployable chatbots with HTML embed code.

\textbf{RAG Integration Fragmentation:} RAG integration varies significantly across frameworks, with each requiring framework-specific configuration. Existing solutions provide workflow-level RAG configuration rather than per-agent document awareness, limiting specialization. RAG systems offer retrieval capabilities but require manual integration with agent systems. AICC-IntelliDoc's DocAware system provides a unified RAG interface with per-agent configuration, supporting multiple search methods through a consistent API.

\textbf{Human-in-the-Loop Pattern Inconsistencies:} HITL support ranges from native first-class features to community workarounds, creating friction when deploying production workflows requiring human oversight. Existing implementations treat HITL as a single interaction point, lacking the distinction between end-user and administrator input modes required for different deployment contexts. AICC-IntelliDoc addresses this with UserProxyAgent supporting both Admin-in-the-Loop and User-in-the-Loop modes.

\section{System Architecture}

AICC-IntelliDoc addresses the fragmentation gap by integrating three core architectural components that work together to enable production-ready chatbot deployment from visual workflows. The framework employs a microservices architecture using PostgreSQL, Milvus (with project-specific collections), Django, and SvelteKit, with complete project-level isolation enabling parallel execution without interference through a five-layer defense-in-depth isolation model: database isolation via foreign key constraints and explicit \texttt{project\_id} filtering, application isolation through \texttt{project.has\_user\_access(user)} validation, vector database isolation with dedicated Milvus collections per project (\texttt{project\_name\_project\_id} pattern), API isolation with project-specific key management and access control, and encryption isolation using PBKDF2 with project-specific keys (100,000 iterations).

The following sections present the three main architectural contributions that address critical gaps identified in Section 2. \textbf{DocAware} (Section 4) addresses the RAG integration fragmentation gap by providing per-agent document awareness with hierarchical document management, enabling agents to access specialized knowledge sections through folder and file-level content filtering. \textbf{Dual-Mode Human-in-the-Loop Architecture} (Section 5) addresses the HITL pattern inconsistency gap by supporting both Admin in the Loop and User in the Loop modes, enabling executive oversight and end-user interactions in different deployment contexts. \textbf{Deployment Architecture} (Section 6) addresses the visual-to-production deployment gap by translating visual workflow designs into deployable chatbots with HTML embed code generation and seamless website integration.

\section{DocAware: Unified RAG System with Hierarchical Document Management}

DocAware is the Retrieval-Augmented Generation (RAG) system that enables agents to retrieve and use document context when generating responses. DocAware provides a unified RAG interface with per-agent configuration and hierarchical document management, enabling specialized knowledge access within the same workflow.

\subsection{Hierarchical Document Structure Management}

The system extracts hierarchical structure from document filenames, enabling folder and file-level content filtering that allows agents to access targeted knowledge sections.

\textbf{Filename-Based Hierarchy Extraction:} During document processing, the system analyzes filenames to extract hierarchical structure (category/subcategory/type). For example, a filename like ``Legal/Contracts/Agreement\_2024.pdf'' is parsed into a virtual folder structure: \texttt{Legal/Contracts/} with the file \texttt{Agreement\_2024.pdf}. This hierarchical path is stored as metadata in the vector database, enabling folder-based filtering.

\textbf{Virtual Folder Structure Creation:} The system builds a virtual folder structure from all document filenames, creating parent folder paths that can be used for filtering. When a document is stored at \texttt{Reports/Financial/Q4\_Report.pdf}, the system creates folder paths for \texttt{Reports/} and \texttt{Reports/Financial/}, enabling filtering at any level of the hierarchy.

\textbf{Hierarchical Path Storage:} Each document chunk in Milvus stores its \texttt{hierarchical\_path} metadata, which contains the full path from root to the document. This enables prefix matching for folder-based filtering: a filter for \texttt{folder\_Reports/Financial} matches all documents whose hierarchical path starts with \texttt{Reports/Financial/}.

\subsection{Per-Agent Configuration with Content Filtering}

Each agent in a workflow can independently configure DocAware with different search methods, parameters, and content filters. This enables specialization: a legal review agent might use hierarchical search with folder-based filtering for contracts, while a technical support agent uses hybrid search with keyword boosting for documentation. DocAware supports multi-select content filtering with OR logic, enabling agents to access specific knowledge sections. Filters can combine multiple folder paths and file IDs, with folder filters using prefix matching and file filters using exact matching. Multiple filters are combined with OR logic, enabling flexible knowledge access patterns.

\subsection{Search Methods}

DocAware integrates with Milvus to support six search strategies: (1) \textit{semantic} search using cosine similarity on embeddings, (2) \textit{hybrid} search combining semantic and keyword matching, (3) \textit{contextual} search using conversation history for query expansion, (4) \textit{hierarchical} search leveraging filename-based structure, (5) \textit{keyword} search with BM25-style ranking, and (6) \textit{threshold-based} search filtering by similarity scores. Each method is accessible through the same unified API, allowing agents to select the optimal strategy for their specialization. All search methods support hierarchical content filtering, enabling targeted retrieval from specific folder structures.

\textbf{Query Extraction and Multi-Input Aggregation:} When agents receive inputs from multiple sources, DocAware intelligently aggregates primary and secondary inputs to construct optimal search queries. The system identifies the primary input (typically from StartNode or previous sequential agent) and secondary inputs (from parallel branches or reflection edges), combining them contextually for retrieval. The query extraction algorithm analyzes conversation history to extract search queries, with optional LLM-based query refinement to preserve concepts and optimize retrieval. The extracted query is then used with the agent's configured search method and content filters to retrieve relevant documents. Document processing details are available in supplementary materials.


\section{Dual-Mode Human-in-the-Loop Architecture}

The UserProxyAgent enables human-in-the-loop interactions, pausing workflow execution to request human input. AICC-IntelliDoc introduces dual input modes supporting different deployment contexts:

\subsection{Admin in the Loop}

Admin in the Loop pauses workflows for administrator input, enabling use cases such as customer service where executives review workflow-generated information and provide expert, monitored responses. Admin in the Loop creates a pause mechanism with different access controls and UI presentation—administrators see the full workflow context, agent reasoning, and retrieved documents before providing curated responses. This enables workflows where automated agents handle routine tasks while human experts provide quality assurance—a pattern critical for compliance-heavy industries (legal, healthcare, finance) where human oversight is mandatory. Admin in the Loop can be implemented in customer service scenarios where administrators leverage domain-specific RAG to gather required information and support end-users with the most relevant answers.

\subsection{User in the Loop}

User in the Loop pauses workflows for end-user input in deployed chatbots, enabling stateful conversations. When a UserProxyAgent in User in the Loop mode is reached, execution pauses and creates a \texttt{HumanInputInteraction} database record with full execution context including input messages from connected agents, conversation history, aggregated input summary, and workflow pause sequence number. The deployed chatbot presents an input interface to the end-user, and upon submission, an API endpoint (\texttt{POST /api/workflow-deploy/<project\_id>/submit-input/}) resumes execution with user input as the node's output. The resume mechanism updates the execution state, records the human response timestamp, and continues workflow execution from the paused state using the human input as the UserProxyAgent's output. UserProxyAgent supports DocAware functionality in User in the Loop mode, allowing user input to trigger document searches and summarization. The ability to switch between Admin in the Loop and User in the Loop gives administrators and users the opportunity to steer agentic workflow outputs toward more desirable results.

\section{Deployment Architecture: Workflow-to-Chatbot Translation}

AICC-IntelliDoc's deployment system translates visual workflows into production-ready chatbots, providing a complete pipeline from visual design to website integration—addressing the visual-to-production deployment gap.

\subsection{Visual Workflow to Chatbot Pipeline}

The deployment pipeline consists of six stages: (1) \textbf{Visual Workflow Design}—users design multi-agent workflows via drag-and-drop interface, placing agent nodes and creating edges representing message flow, (2) \textbf{JSON Storage and Validation}—workflow graphs are stored as JSON and validated for cycles and required StartNode, (3) \textbf{Execution Engine Integration}—validated workflows are integrated with the orchestration engine for processing user queries, (4) \textbf{Deployment Configuration}—users select a workflow, configure allowed origins, set per-origin rate limits, and specify initial greeting, (5) \textbf{HTML Embed Code Generation}—the system generates self-contained HTML code with chatbot interface, API endpoint URLs, and session management, and (6) \textbf{Website Integration}—users copy the generated HTML code and paste it into their websites, with the deployment system validating allowed origins and enforcing rate limits.

\subsection{HTML Embed Code Generation and Session Management}

The deployment system generates HTML embed code that includes a complete chatbot interface with responsive design, message display area with user and assistant message styling, input interface for user queries, API endpoint integration for workflow execution, session management for stateful conversations, and error handling with loading states. The generated code is self-contained and requires only the endpoint URL and initial greeting configuration. Each chatbot session is assigned a unique session ID, and conversation history is stored as JSON arrays containing message sequences, enabling chatbots to maintain context across multiple user interactions.

\subsection{Allowed Origins and Rate Limiting}

The deployment system manages a list of allowed origins (domains) that can embed and access the chatbot. Each deployment can have multiple allowed origins, enabling the same workflow to be embedded in multiple websites. The system validates the \texttt{Origin} header in incoming requests, ensuring that only requests from allowed origins are processed. Each allowed origin can have its own rate limit configuration (requests per minute), enabling fine-grained control over chatbot usage. The system tracks request counts per origin using in-memory counters and enforces rate limits before processing workflow execution requests. CORS middleware handles cross-origin requests for chatbot embedding, validating origin headers, setting appropriate CORS response headers, and handling preflight OPTIONS requests. Each deployment is associated with a specific workflow, enabling multiple deployments per project with different workflows and active/inactive toggles for enabling or disabling chatbot access.

\section{Limitations, Ethical Considerations, and Conclusion}

The current work presents system architecture and implementation with initial validation tests. Key limitations include: empirical evaluation is limited to isolation validation; user studies validating visual workflow designer accessibility claims are not yet conducted; and security analysis of isolation mechanisms is informal. AICC-IntelliDoc handles sensitive documents requiring strong privacy protections. The isolation architecture provides encryption at rest and project-specific API key encryption, though key rotation mechanisms and formal security analysis are needed. Documents sent to LLM APIs may be logged by providers, requiring user awareness. The system does not filter harmful or biased content, which may influence AI-generated summaries and agent responses. Responsible use policies are encouraged.

This paper presents AICC-IntelliDoc, a framework that translates visual AI agentic workflows into production-ready chatbots. The framework addresses a critical gap where existing solutions require extensive integration or lack deployment capabilities. The integration of dual-mode human-in-the-loop architecture, hierarchical document management, and workflow-to-chatbot translation creates a practical solution for building document-intelligent conversational applications. AICC-IntelliDoc will be released as open source, with complete source code, documentation, and deployment configurations made publicly available.

\section*{Acknowledgments}

We thank the open-source community for the excellent tools and frameworks that made this work possible.

% Entries for the entire Anthology, followed by custom entries
% Bibliography entries are provided in custom.bib
% Note: Some entries in custom.bib may need verification and updates with exact publication details
\bibliography{anthology,custom}
\bibliographystyle{acl_natbib}

\end{document}

